% Обсуждение и ограничения (main.tex)
\section{Обсуждение и ограничения}\label{sec:discussion}

\paragraph{Связь с дискуссией о <<чистоте>> сетей.}
Замена слоя явным алгоритмом с малым KL "--- конструктивная форма
утверждения <<мы понимаем, что делает этот слой>> в смысле критерия
Gao~\cite{gao_notclean}. По этому критерию первые слои Qwen3 <<чисты>>
наполовину: их внимание действительно локально и сводится к окну с
потерей 1.5--1.8\,\% PPL, а их MLP "--- нет: лучшая явная замена с
предсказанием нейронов стоит 4.7\,\% на 0.6B и 1.7\,\% на 1.7B, табличная
замена по токену "--- 12\,\%, линейная "--- 20\,\%. При этом MLP поздних
слоёв по карте заменяемости в 15--20 раз <<чище>> ранних. Иными словами,
простота первых слоёв "--- свойство их внимания, а не их MLP.

\paragraph{Четыре практических правила.} Из результатов следуют правила
для любой пост-хок замены блоков в замкнутой форме: (1) суммы по токенам
не должны включать строки с массивной активацией "--- одна такая строка
в калибровке доминирует в задаче наименьших квадратов, а порог опасности
лежит между одним и четырьмя окнами из 256; (2) сток "--- это массивная активация, где бы она ни возникла:
нейроны, создающие её, должны сохраняться явно, а окно внимания "---
хранить несущие её ключи; у Qwen3 это позиция~0, у Llama-семейства "---
ещё и токены вроде <<the>> в любой позиции, и без этого правила окно
стоит 64\,\% вместо 4\,\% на SmolLM2; (3) выбирать блок для замены
надо по измеренной чувствительности, а не по косинусной близости
состояний: на Qwen3-0.6B косинусный критерий ошибается в 4--5 раз;
(4) калибровать надо на смеси доменов "--- иначе ошибка калиброванной
замены вне домена втрое выше.

\paragraph{Ограничения.}
\begin{itemize}
  \item Ускорение оценивается по числу операций. Измеренное время
    (табл.~\ref{tab:timing}) показывает, что без специализированного
    ядра окна реализация на чистом PyTorch медленнее fused-ядра полного
    внимания; экономия в несколько процентов умножений на первых слоях
    сама по себе ускорения не даёт, выигрыш окна "--- память кэша
    ключей-значений на длинном контексте.
  \item Калибровка на одном домене переобучает калиброванные замены MLP
    (ошибка вне домена втрое выше); смешанная калибровка выравнивает
    ошибку ценой роста внутри домена.
  \item Исследованы два семейства (Qwen3-0.6B/1.7B, SmolLM2-1.7B,
    Llama-3.2-1B) размером до 1.7B; вычисления выполнены на ноутбуке.
    Выводы о внимании и о стоках повторяются на обоих семействах, выводы о
    заменимости MLP и об априорном критерии для MLP "--- нет: у
    Llama-семейства первый MLP дороже, а концентрация вклада предсказывает
    KL слабее. Вопрос о сохранении <<обменного курса>> с ростом размера
    требует моделей 4B--32B. Кривые по $k$ для Llama-3.2-1B и точки окна с
    массивными стоками при всех $k$ досчитываются; в таблицах они
    отмечены.
  \item Бенчмарки посчитаны на подвыборках: по 200 примеров для всех
    столбцов ($\pm2.5$ п.п. для среднего по семи задачам) и по 1000 для
    главных ($\pm1.1$ п.п.); интервалы для PPL и KL даны по бутстрэпу окон.
  \item Дистилляционный контроль (300 шагов AdamW) ограничен шестью
    методами и малым числом шагов; для метода O4 он расходится.
  \item Методы из литературы адаптированы к Qwen3 (A3) или взяты в
    детерминированной части (SVD-LLM без дообучения, Deja Vu в
    оракульном режиме, MoEfication без обучаемого маршрутизатора). Их
    обучаемые варианты оценены только в дистилляционном контроле.
\end{itemize}

\section{Заключение}
Мы спросили, можно ли заменить первые слои готовой LLM явными
детерминированными алгоритмами, и получили ответ в четырёх частях.
Внимание первых слоёв "--- да: окно из 64 позиций со стоком заменяет
первые три слоя Qwen3-0.6B и 1.7B с ростом PPL 1.5--1.8\,\% и без изменения
точности на бенчмарках, первые шесть "--- в пределах 5\,\%, и число стоков,
как и длина контекста до 8 тыс. токенов, на это почти не влияют. MLP
первых слоёв "--- лишь частично: предсказание 25\,\% активных нейронов в
замкнутой форме (M-O4) стоит 4.7 и 1.7\,\% PPL при стоимости $0.3$ MLP,
тогда как пороговый отбор, линейные карты и низкоранговые приближения
теряют в разы больше, а по карте заменяемости MLP ранних слоёв "--- самые
незаменимые во всей модели. Любая калибровка в замкнутой форме "--- только
с защитой от массивных активаций: одна строка данных отравляет регрессию
(что мы доказали формулой, воспроизвели дозой и вылечили исключением
строк), а один нейрон, создающий сток, должен быть сохранён. Второе
семейство показало общий вид этого правила: сток "--- массивная
активация, у SmolLM2 её создаёт один нейрон на токене <<the>> в любой
позиции, и только защита таких нейронов и хранение таких ключей делают
замены переносимыми (окно: 64 $\to$ 4\,\%; M-O4: 123 $\to$ 11\,\%). Наконец,
дешёвые априорные статистики предсказывают заменяемость слоёв и на
непросмотренной модели выбирают план замен в 10--15 раз лучше префикса.
Все замены укладываются в один обменный курс "--- около процента PPL за
процент вычислений на 0.6B и вдвое лучше на 1.7B, "--- и проверка того,
продолжает ли курс улучшаться с размером, составляет следующий шаг.
